% BOOK-FILE-ID: f-3076141c1a19
% subsubsection: 一次分数型  \, $\displaystyle a_{n+1}=\frac{pa_n+q}{ra_n+s}$

\subsubsection{\texorpdfstring{一次分数型   $a_{n+1}=\frac{pa_n+q}{ra_n+s}$}{一次分数型}}\label{節：一次分数型}
\booknavref{節：逆数型}{逆数型}の一般形である,
\[
  a_{n+1}=\frac{pa_n+q}{ra_n+s}
\]
で表される漸化式について考えよう.

以下では$r\ne0$, $ps-qr\ne0$とし,元の分母が$0$にならない範囲で議論する.\footnote{$r=0,\ s\ne0$なら特性方程式型へ戻る.
  $r\ne0,\ ps-qr=0$なら,分母が0でない入力に対して右辺は定数$\frac{p}{r}$になる.次の分母も確認する.例えば$\frac{a_n-1}{a_n-1}$は,初項が$1$でなくても次項が$1$になり,その次の項は定義できない.}
  
既に学んだ逆数型に帰着させることを目標に式変形を進める.
つまり,
\[
  a_{n+1}=\frac{pa_n+q}{ra_n+s} \iff b_{n+1}=\frac{tb_n}{ub_n+v}
\]
となるような$b_n$と$t,u,v$の組を探せば良い.
言い換えれば$q$を消すということになるので,特性方程式型の考え方を用いて,ある不動点を見つけ,その値を引くことを目標にしよう.

$b_n=a_n-x$と置き,分子の定数項が消えるように$x$を選ぶ.
$a_n=b_n+x$を代入すると,
\BookMathLayout{\begin{align*}
    b_{n+1}&=\frac{p(b_n+x)+q}{r(b_n+x)+s}-x\\
    &=\frac{(p-rx)b_n+px+q-rx^2-sx}
    {rb_n+rx+s}
  \end{align*}}{\[
    \begin{aligned}
      &b_{n+1}\\
      &{}=\frac{p(b_n+x)+q}{r(b_n+x)+s}-x\\
      &{}=\frac{(p-rx)b_n+px+q-rx^2-sx}
      {rb_n+rx+s}
    \end{aligned}
  \]}
従って,
\[
  rx^2+(s-p)x-q=0
\]
となれば$q$が消えて$b_n$と$t,u,v$が得られる.
また,この二次方程式も特性方程式の一種である.
\[
  rx^2+(s-p)x-q=0
  \quad\Longleftrightarrow\quad
  x=\frac{px+q}{rx+s}
\]
であるから,形式的には$a_{n+1}$や$a_n$を$x$に置き換えた式について,これを満たす$x$,すなわち元の\bookterm{fixed}{不動点}を引けば,
\[
  b_{n+1}=\frac{(p-rx)b_n}{rb_n+rx+s}
\]
という逆数型になる.ここで,\,$p-rx$も$rx+s$も0でない.\footnote{特性方程式より,\,$(p-rx)(rx+s)=ps-qr\ne0$.}
初項が$x$なら定数数列である.それ以外では$b_n$は途中で0にならず,逆数を取って解ける.

ここまでで逆数型として解くことができるが,定数項を消せる値を求める問題は,
\[
  rx^2+(s-p)x-q=0
\]
という二次方程式を解く問題に帰着した.
この方程式が異なる2つの実数解$\alpha,\beta$を持つなら,その左辺は
\[
  rx^2+(s-p)x-q=r(x-\alpha)(x-\beta)
\]
と因数分解できる.
$\alpha,\beta$は別々に選ぶ定数ではなく,もとの漸化式の係数から決まる,同じ二次方程式の二つの解なのである.
そして,どちらの解も定数項を消す条件を満たすので,\, $a_n-\alpha$と$a_n-\beta$のどちらについても,同じ変形ができる.
先ほどの逆数型の式を,もとの$a_n$を使って書き直すと,
\begin{align*}
  a_{n+1}-\alpha
  &=\frac{p-r\alpha}{ra_n+s}(a_n-\alpha),\\
  a_{n+1}-\beta
  &=\frac{p-r\beta}{ra_n+s}(a_n-\beta)
\end{align*}
となる.

この二つの式から,さらに単純な法則を取り出せないだろうか.
それぞれの差に掛かる係数には$a_n$が含まれているので,差そのものは一般には等比数列にならない.
しかし,2つの係数には同じ分母$ra_n+s$が現れている.
そこで2つの式の比を取れば,この分母が消え,定数の係数だけが残りそうである.

初項が$\alpha$または$\beta$なら定数数列になるので,先に扱っておこう.
以下では初項がどちらの不動点とも異なる場合を考える.
また,\, $p-r\beta=0$なら不動点の条件から$q=s\beta$となり,\, $ps-qr=0$となって仮定に反する.
よって,\, $p-r\beta\ne0$であり,上の式から$a_n\ne\beta$なら$a_{n+1}\ne\beta$である.
初項から順に考えると,すべての項で$a_n\ne\beta$なので,二つの差の比を取ることができる.
実際に,
\begin{align*}
  \frac{a_{n+1}-\alpha}{a_{n+1}-\beta}
  &=\cfrac{
    \cfrac{(p-r\alpha)(a_n-\alpha)}{ra_n+s}}
    {
    \cfrac{(p-r\beta)(a_n-\beta)}{ra_n+s}}\\
  &=\frac{p-r\alpha}{p-r\beta}
    \frac{a_n-\alpha}{a_n-\beta}
\end{align*}
となる.
したがって,2つの差の比は,公比$\frac{p-r\alpha}{p-r\beta}$の等比数列となる.
分子の定数項を消せる二つの式を比べることで,等比数列の法則を見つけられたのである.

ここで
\[
  b_n=\frac{a_n-\alpha}{a_n-\beta}
\]
と置けば,
\[
  b_{n+1}
  =
  \frac{p-r\alpha}{p-r\beta}b_n
\]
となる.したがって,\,$b_n$は等比数列である.
$\alpha$と$\beta$の大小関係は,解の大小関係を要求されていないことから自由に選んでよい.
つまり,\,一次分数型漸化式では,
\[
  a_{n+1}=\frac{pa_n+q}{ra_n+s}
\]
という複雑な形をそのまま扱うのではなく,
\[
  \frac{a_n-\alpha}{a_n-\beta}
\]
という「$\alpha,\beta$からの差の比」に置き換えることで,
等比数列に変えることができる.

そして,\,$\alpha,\beta$は,
\[
  rx^2+(s-p)x-q=0
\]
の異なる2解として自然に現れる.

したがって\,$\alpha,\beta$は,もとの漸化式の$a_n$を形式的に$\alpha,\beta$に置き換えたときにも値が変化しない数,すなわち不動点である.

なお,\,$rx^2+(s-p)x-q=0$が重解$\alpha$を持つ場合は,\,$\beta$と異なる2つの不動点を取れないため,上の議論はそのまま使えない.
このときは$\alpha=\frac{p-s}{2r}$,すなわち$s=p-2r\alpha$が成り立つため,
\[
  c=p-r\alpha=r\alpha+s\ne0
\]
とおくと,
\[
  a_{n+1}-\alpha
  =\frac{c(a_n-\alpha)}{r(a_n-\alpha)+c}
\]
となる.
初項が$\alpha$なら定数数列になるので,以下では$a_1\ne\alpha$とする.
上の式では$c\ne0$であり,元の分母が0でない範囲では$a_n-\alpha$は途中で0にならない.
\[
  b_n=a_n-\alpha
\]
とおけば逆数を取ることができ,
\[
  \frac1{b_{n+1}}=\frac1{b_n}+\frac rc
\]
という等差数列が得られ,ここから$a_n$を求めることができる(具体的な計算は例題(3)を参照).

\begin{bookpoint}
  \textbf{POINT}\par
  $a_{n+1}$と$a_n$を$x$に置き換えた二次方程式を解き,異なる2解$\alpha,\beta$を得たとき,\,$b_n=\frac{a_n-\alpha}{a_n-\beta}$とおく.
  初項が不動点のときは定数数列として分け,分母が0でないことを確かめる.
  $b_n$について等比型漸化式が作れるので,その一般項を得たら$a_n$に戻す.
\end{bookpoint}

\noindent \bookblocklink{example-mobius-1}{problem}{answer}{【例題】}\par\nobreak
\begin{enumerate}[label=(\arabic*),parsep=3pt,format=\bookitemlink{example-mobius-1}{problem}{answer}]
  \item $a_1=3,\, a_{n+1}=\frac{2a_n+1}{a_n+2}$
  \item $a_1=0,\, a_{n+1}=\frac{a_n+4}{a_n+1}$
  \item $a_1=5,\, a_{n+1}=\frac{4a_{n}-9}{a_{n}-2}$
  \item $a_1=1,\, a_{n+1}=1+\frac{1}{a_n}$
\end{enumerate}

\noindent\bookblocklink{example-mobius-1}{hint}{problem}{【方針】}\par\nobreak
\begin{enumerate}[label=(\arabic*),parsep=3pt,format=\bookitemlink{example-mobius-1}{hint}{problem}]
  \item 不動点$1,-1$からの差の比を作ると,共通の分母$a_n+2$が消えます.
  \item 同様に\bookterm{fixed}{不動点}$\alpha,\beta$を求めましょう.今回は差の比の公比が負の値になります.
  \item 今回は\bookterm{fixed}{不動点}の方程式が重解になります.\,$\frac{a_n-\alpha}{a_n-\beta}$は使えないので,\,$b_n=a_n-\alpha$と置いてから逆数型に帰着させる別の手法を使います.
  \item このような形も一次分数型です.解と係数の関係を活用して解きましょう.
\end{enumerate}

\noindent\bookblocklink{example-mobius-1}{answer}{problem}{【解答】}\par\nobreak
\begin{enumerate}[label=(\arabic*),parsep=3pt,format=\bookitemlink{example-mobius-1}{answer}{problem}]
  \item $a_n=\frac{2\cdot3^{n-1}+1}{2\cdot3^{n-1}-1}$
  \item $a_n=\frac{2\left(3^{n-1}+(-1)^n\right)}{3^{n-1}-(-1)^n}$
  \item $a_{n}=\frac{6n-1}{2n-1}$
  \item $a_n=\cfrac{\left(\cfrac{1-\sqrt{5}}{2}\right)^{n+1}-\left(\cfrac{1+\sqrt{5}}{2}\right)^{n+1}}{\left(\cfrac{1-\sqrt{5}}{2}\right)^n-\left(\cfrac{1+\sqrt{5}}{2}\right)^n}$
\end{enumerate}

\noindent\bookblocklink{example-mobius-1}{solution}{problem}{【解法】}\par\nobreak
\begin{enumerate}[label=(\arabic*),parsep=3pt,format=\bookitemlink{example-mobius-1}{solution}{problem}]
  \item $a_1=3,\, a_{n+1}=\frac{2a_n+1}{a_n+2}$\\
    \bookterm{fixed}{不動点}の方程式$x=\frac{2x+1}{x+2}$の解は$1,-1$である.
    $b_n=\frac{a_n-1}{a_n+1}$と置くと,
    \begin{align*}
      b_{n+1}&=\cfrac{\cfrac{2a_n+1}{a_n+2}-1}{\cfrac{2a_n+1}{a_n+2}+1}\\
      &=\frac{2a_n+1-(a_n+2)}{2a_n+1+(a_n+2)}\\
      &=\frac13\frac{a_n-1}{a_n+1}=\frac13b_n
    \end{align*}
    $b_1=\frac{3-1}{3+1}=\frac{1}{2}$なので,
    \[
      b_n=\frac12\left(\frac13\right)^{n-1}
      =\frac1{2\cdot3^{n-1}}
    \]
    である.元の項への復元は,
    \begin{align*}
      b_n(a_n+1)&=a_n-1,\\
      a_n&=\frac{1+b_n}{1-b_n}
      =\frac{2\cdot3^{n-1}+1}{2\cdot3^{n-1}-1}
    \end{align*}
    となる.
  \item $a_1=0,\, a_{n+1}=\frac{a_n+4}{a_n+1}$\\
    \bookterm{fixed}{不動点}の方程式$x=\frac{x+4}{x+1}$の解は$2,-2$である.
    公比は$\frac{1-2}{1+2}=-\frac{1}{3}$なので,\,$b_n=\frac{a_n-2}{a_n+2}$と置くと,
    \begin{align*}
      b_{n+1}&=-\frac13b_n,\quad b_1=-1,\\
      b_n&=(-1)\left(-\frac13\right)^{n-1}
      =\frac{(-1)^n}{3^{n-1}}
    \end{align*}
    を得る.\,$b_n(a_n+2)=a_n-2$から復元して,
    \begin{align*}
      a_n&=\frac{2(1+b_n)}{1-b_n}\\
      &=\frac{2\left(3^{n-1}+(-1)^n\right)}{3^{n-1}-(-1)^n}
    \end{align*}
    となる.
  \item $a_1=5,\, a_{n+1}=\frac{4a_n-9}{a_n-2}$\\
    \bookterm{fixed}{不動点}の方程式$x=\frac{4x-9}{x-2}$は重解$3$を持つ.
    $b_n=a_n-3$と置くと,
    \begin{align*}
      b_{n+1}&=\frac{4a_n-9}{a_n-2}-3
      =\frac{a_n-3}{a_n-2}\\
      &=\frac{b_n}{b_n+1},\quad b_1=2
    \end{align*}
    となる.正の$b_n$が続くので,\,$b_n$も元の分母$b_n+1$も0でない.
    逆数を取れば,
    \begin{align*}
      \frac1{b_{n+1}}&=\frac1{b_n}+1,\\
      \frac1{b_n}&=\frac12+(n-1)=\frac{2n-1}{2}
    \end{align*}
    である.従って,
    \[
      a_n=3+b_n=3+\frac2{2n-1}=\frac{6n-1}{2n-1}
    \]
    を得る.
  \item $a_1=1,\, a_{n+1}=1+\frac1{a_n}$\\
    右辺を$\frac{a_n+1}{a_n}$と書けば,一次分数型である.
    \bookterm{fixed}{不動点}の方程式
    \[
      x=\frac{x+1}{x}\booktag{mobius-example-fixed}{1}
    \]
    の解を
    \[
      \alpha=\frac{1+\sqrt5}{2},\qquad
      \beta=\frac{1-\sqrt5}{2}
    \]
    とする.\bookeqbackref{mobius-example-fixed}{(1)}の解と係数の関係から,
    \[
      \alpha+\beta=1,\qquad \alpha\beta=-1
      \booktag{mobius-example-roots}{2}
    \]
    である.\,$b_n=\frac{a_n-\alpha}{a_n-\beta}$と置くと,
    \begin{align*}
      b_{n+1}&=\frac{a_n+1-\alpha a_n}{a_n+1-\beta a_n}\\
      &=\frac{(1-\alpha)a_n+1}{(1-\beta)a_n+1}
      \booktag{mobius-example-ratio}{3}
    \end{align*}
    となる.\bookeqbackref{mobius-example-roots}{(2)}を\bookeqbackref{mobius-example-ratio}{(3)}に使えば,
    \begin{align*}
      b_{n+1}&=\frac{\beta a_n-\alpha\beta}{\alpha a_n-\alpha\beta}\\
      &=\frac{\beta(a_n-\alpha)}{\alpha(a_n-\beta)}
      =\frac\beta\alpha b_n
    \end{align*}
    である.\,$b_1=\frac{1-\alpha}{1-\beta}=\frac{\beta}{\alpha}$より,
    \[
      b_n=\left(\frac\beta\alpha\right)^n
    \]
    となる.復元式$a_n=\frac{\beta b_n-\alpha}{b_n-1}$に代入し,分母分子に$\alpha^n$を掛けると,
    \begin{align*}
      a_n&=\frac{\beta^{n+1}-\alpha^{n+1}}{\beta^n-\alpha^n}\\
      &=\cfrac{\left(\cfrac{1-\sqrt5}{2}\right)^{n+1}-\left(\cfrac{1+\sqrt5}{2}\right)^{n+1}}
      {\left(\cfrac{1-\sqrt5}{2}\right)^n-\left(\cfrac{1+\sqrt5}{2}\right)^n}
    \end{align*}
    を得る.
\end{enumerate}
\paragraphbf{特性方程式が重解の時}
  特性方程式が重解$\alpha$を持つときは,差の比を作れない.
  先ほど見たように,\,$a_1=\alpha$を定数数列として分け,それ以外では$c=p-r\alpha$とすると,
  \[
    \frac1{a_{n+1}-\alpha}
    =\frac1{a_n-\alpha}+\frac rc
  \]
  という\bookterm{arithmetic}{等差数列}に帰着する.
  \textbf{異なる不動点なら二つの差の比,重解なら一つの差の逆数}を$b_n$と置けばよい.
  例題(3)は$\alpha=3,\ r=c=1$なので,逆数が毎回$1$ずつ増えた.

\paragraphbf{小話：フィボナッチ数列}
  例題(4)は,フィボナッチ\bookterm{sequence}{数列}$F_1=F_2=1,\ F_{n+2}=F_{n+1}+F_n$と関係している.
  隣接項の比$\frac{F_{n+1}}{F_n}$は同じ漸化式$a_{n+1}=1+\frac{1}{a_n}$と初項$1$を満たすので,\,$a_n=\frac{F_{n+1}}{F_n}$である.
  この比の\bookterm{limit}{極限}は黄金比$\frac{1+\sqrt5}{2}$となる.\footnote{リュカ数列も同じ\bookterm{characteristic}{特性方程式}$x^2=x+1$を持ち,隣接項の比が黄金比へ収束する.
    トリボナッチ数列では$x^3=x^2+x+1$となり,比の極限は約$1.839$のトリボナッチ定数になる.}

\paragraphbf{数学的帰納法による一般項の証明}
  一次分数型では,初めの数項から一般項を予想し,帰納法で確かめる方法もある.
  数項を求める誘導があれば,分母と分子の増え方にも注目しよう.例題(3)で試す.\\
  \noindent\bookblocklink{example-mobius-2}{problem}{answer}{【例題】}\par\nobreak
  \begin{quote}
    以下の\bookterm{recurrence}{漸化式}
    \[
      a_1=5,\, a_{n+1}=\frac{4a_{n}-9}{a_{n}-2}
    \]
    を満たす\bookterm{sequence}{数列}$a_n$の\bookterm{general}{一般項}を帰納法を用いて求めよ.
  \end{quote}
  \bookblocklink{example-mobius-2}{hint}{problem}{【方針】}
  \begin{quote}
    実験的に代入計算をしてみると,
    \[
      a_1=\frac{5}{1},\,a_2=\frac{11}{3},\,a_3=\frac{17}{5},\,a_4=\frac{23}{7},\,a_5=\frac{29}{9}
    \]
    となる.分母が2ずつ,分子が6ずつ増えていることから,\bookterm{general}{一般項}は
    \[
      a_n=\frac{6n-1}{2n-1}
    \]
    と予想する.
  \end{quote}
  \bookblocklink{example-mobius-2}{answer}{problem}{【解答】}
  \begin{quote}
    予想した\bookterm{general}{一般項}を確かめる.\,$n=1$では$\frac{6-1}{2-1}=5$なので成立する.
    正の整数$k$について,
    \[
      a_k=\frac{6k-1}{2k-1}
    \]
    が成り立つと仮定する.このとき,\bookterm{recurrence}{漸化式}より,
    \begin{align*}
      a_{k+1}&=\frac{4(6k-1)-9(2k-1)}{6k-1-2(2k-1)}\\
      &=\frac{6k+5}{2k+1}
      =\frac{6(k+1)-1}{2(k+1)-1}
    \end{align*}
    分母$a_k-2=\frac{2k+1}{2k-1}$も0でなく,\,$n=k+1$で成立する.
    帰納法より,\,$a_n=\frac{6n-1}{2n-1}$がすべての正の整数$n$で成り立つ.\qed

  \end{quote}
